Pretraining a motion representation on abundant vision-captured skeleton data and transferring it to data-scarce body-worn inertial sensors is an attractive strategy for wearable gesture recognition, but \textit{how} the representation should be pretrained --- with a supervised, reconstruction, contrastive, or hybrid objective --- is unsettled, and the field has drifted toward reconstruction/self-supervised objectives on the assumption that they are the uniformly safe choice under domain shift. We show instead that the \textit{identity} of the unsafe objective is not a fixed property of the objective --- it depends on how wide the gap is. We present a controlled study isolating the pretraining objective across five variants (scratch, supervised, reconstruction, contrastive, hybrid) while transferring an NTU RGB+D Transformer encoder to a small Xsens gesture dataset (22 gestures, 5 subjects) over a shared, skeleton-derived quaternion representation --- a moderate device gap --- under subject-held-out few-shot calibration. Objective utility is gap-gated: a 45-point within-domain accuracy spread compresses to a few points here. At this moderate gap, and against the prevailing assumption, the significant negative-transfer case is \textit{pure masked reconstruction}, which transfers below random initialization at every shot count (clip-level McNemar p<0.001, three seeds), while the hybrid objective is robustly best and the label-supervised objectives --- softmax and contrastive alike --- are competitive and best-calibrated. Extending to two independent public datasets and a substantially wider gap --- a true cross-modal wearable-IMU target with no skeleton-derived representation to fall back on --- reverses this: the label-supervised objectives, dominant at small gaps, become the \textit{worst} priors, while the hybrid remains safe. The negative-transfer culprit thus flips identity with gap width; a reconstruction component is the one ingredient that stays safe across both regimes tested. Layer-wise CKA shows representational alignment is necessary but not sufficient in either regime. The prior's more gap-stable value is calibration and data-efficiency: it converges faster and helps most (+27 points) with 0--1 enrolled subjects, washing out by four --- though this cold-start advantage also disappears on a target with a strong representation of its own. An abstract controller driven by the recognizer's own posteriors shows compressed recognition differences compound into much larger task-level ones: across 120 randomly-drawn task designs, two outcome models, and a tuning-free safety operating point, the reconstruction objective compounds worst under every stress test, while the best-calibrated objective yields the safest confidence-threshold operating points.
